% Work/MAT – 1. ročník – hodina 51. Navazuje na 49–50.
% Primární metodický zdroj: Realisticky.cz, 010809 Rozklad mnohočlenů na součin IV,
% uživatelem dodaný dokument „09 Rozklad mnohočlenů na součin IV.pdf“, zejména Př. 1–5.
% Výběr a zápis příkladů včetně řešení autorsky přepracován; nejde o opis lekce.
\documentclass[11pt,a4paper]{article}
% Spolecny styl pro ucitelske pripravy MAT (Work/MAT)
% Verze 1.1, 2026-10-23
\NeedsTeXFormat{LaTeX2e}
\ProvidesPackage{priprava}[2026/10/23 v1.1 Sablona ucitelskych priprav MAT]

% Zakladni typografie a cestina
\RequirePackage[T1]{fontenc}
\RequirePackage[utf8]{inputenc}
\RequirePackage[czech]{babel}
\RequirePackage{lmodern}
\RequirePackage{microtype}
\RequirePackage{amsmath,amssymb}
\RequirePackage{xcolor}
\RequirePackage{enumitem}
\RequirePackage[most]{tcolorbox}
\RequirePackage{tabularx,booktabs,array}
\RequirePackage{fancyhdr}
\RequirePackage{needspace}
\RequirePackage{geometry}
\RequirePackage[hidelinks]{hyperref}

\geometry{a4paper,top=17mm,bottom=18mm,left=23mm,right=23mm}
\definecolor{notegray}{gray}{0.39}
\setlength{\parindent}{0pt}
\setlength{\parskip}{2.5pt plus 1pt minus 0.5pt}
\setlist[itemize]{topsep=2pt,itemsep=1pt,parsep=0pt,leftmargin=1.45em}
\setlist[enumerate]{topsep=2pt,itemsep=2pt,parsep=0pt,leftmargin=1.65em}
\setcounter{secnumdepth}{0}
\renewcommand{\arraystretch}{1.12}

% Zahlavi: vlevo rocnik | globalni tema, vpravo cislo: nazev hodiny
% Zapati zustava zcela prazdne.
\pagestyle{fancy}
\fancyhf{}
\renewcommand{\headrulewidth}{0.35pt}
\renewcommand{\footrulewidth}{0pt}
\setlength{\headheight}{14pt}
\newcommand{\nastavPripravu}[4]{%
  \fancyhead[L]{\small #1\enspace\textbar\enspace #2}%
  \fancyhead[R]{\small #3: #4}%
}

% Minimalisticke cernobile boxy. Kazda definice zvlast.
\tcbset{commonlesson/.style={enhanced,breakable=false,colback=white,colframe=black,
  colbacktitle=white,coltitle=black,boxrule=0.45pt,
  arc=0.8mm,outer arc=0.8mm,left=3mm,right=3mm,top=1.6mm,bottom=1.6mm,
  toptitle=1mm,bottomtitle=0.7mm,fonttitle=\bfseries,
  attach boxed title to top left={xshift=2.5mm,yshift=-1.15mm},
  boxed title style={colback=white,colframe=white,boxrule=0pt,left=1mm,right=1mm}}}
\newtcolorbox{zadaniBox}[1][]{commonlesson,title={Zadání},#1}
\newtcolorbox{prikladBox}[1][]{commonlesson,title={Příklad},#1}
\newtcolorbox{definiceBox}[1][]{commonlesson,title={Definice},#1}

% Spolecny box pro zadani a jeho reseni, bez mezery mezi castmi.
% Pouziti:
%   \begin{prikladReseniBox}
%     Zadani ulohy.
%     \tcblower
%     \begin{align*} ... \end{align*}
%   \end{prikladReseniBox}
% Spodni cast ma jemny sedy podklad a delici linku; lze v ni sazet
% vice radku matematiky. Obe casti maji jen jeden spolecny ramecek.
\tcbset{solutionlesson/.style={
  commonlesson,
  enhanced jigsaw,
  breakable,
  colbacklower=black!3,
  segmentation style={solid,draw=black!25,line width=0.35pt},
  fontlower=\small\color{notegray},
  before lower={\noindent\textbf{Řešení:}\enspace},
  leftlower=3mm,rightlower=3mm
}}
\newtcolorbox{prikladReseniBox}[1][]{solutionlesson,title={Příklad},#1}
\newtcolorbox{zadaniReseniBox}[1][]{solutionlesson,title={Zadání},#1}

% Pomucky pro autora pripravy: sedive poznamky; \reseni pro starsi pripravy
\newcommand{\poznamka}[1]{\par\begingroup\small\color{notegray}\itshape
  \textnormal{\textbf{Poznámka pro učitele.}}\enspace #1\par\endgroup}
\newcommand{\reseni}[1]{\par\begingroup\small\color{notegray}
  \textbf{Očekávané řešení:}\enspace #1\par\endgroup}
\newcommand{\blok}[1]{\needspace{5\baselineskip}\subsection*{#1}}
\newcommand{\vyrok}[1]{\uv{#1}}

\nastavPripravu{1. ročník}{Algebraické výrazy}{51}{Rozklady rozdělením a seskupením členů}
\begin{document}
\poznamka{45minutové jádro: 5 min opakování vytýkání a trojčlenů, 10 min odvození principu (P1), 10 min samostatně samostatná práce a kontrola, 8 min složitější seskupování (P2), 9 min druhá samostatná práce a kontrola, 3 min shrnutí. poslední samostatná úloha a rozšíření jsou rezervy. Domácí úlohy 15–20 minut. Vysvětli, že jeden člen lze rozdělit více způsoby, ale hledáme takový, který vytvoří společnou závorku.}
\blok{Když nestačí vytýkat: vytvořme vhodné skupiny}
V hodině 49 jsme vytýkali ze čtyř členů. Nyní máme často pouze tři členy. Jeden z nich proto můžeme \emph{rozdělit} na dva sčítance a získané čtyři členy vhodně seskupit. Nejde o nový vzorec: všechny kroky vycházejí z distributivity.
\begin{prikladReseniBox}[breakable=false,title={Příklad 1: jeden člen nahradíme součtem dvou}]
Rozlož $2x^2+7x+3$ na součin seskupováním.
\tcblower
Hledáme rozdělení prostředního členu $7x=6x+x$:
\begin{align*}
2x^2+7x+3&=2x^2+6x+x+3\\
&=2x(x+3)+(x+3)\\
&=(2x+1)(x+3).
\end{align*}
Zkouška: $(2x+1)(x+3)=2x^2+7x+3$.
\end{prikladReseniBox}
\poznamka{Rozdělení $7x=6x+x$ není nová definice ani zvláštní pravidlo.
Je to vhodná úprava zápisu podle distributivity. Ptej se, proč právě tato dvojice
umožní získat dvakrát stejnou závorku.}
\poznamka{Nejde o automatický algoritmus pro každý trojčlen. U $ax^2+bx+c$ s celočíselnými koeficienty často pomůže najít dvě čísla se součtem $b$ a součinem $ac$; pak rozdělíme člen $bx$. Tento tip zaváděj až po P1, ne jako náhradu za zdůvodnění.}
\begin{prikladReseniBox}[breakable=false,title={Příklad 2: střídání znamének}]
Rozlož $3x^2-5x-2$.
\tcblower
Rozdělíme $-5x=-6x+x$:
\begin{align*}
3x^2-5x-2&=3x^2-6x+x-2\\
&=3x(x-2)+(x-2)\\
&=(3x+1)(x-2).
\end{align*}
\end{prikladReseniBox}
\begin{zadaniReseniBox}[breakable=false,title={Samostatná práce: rozděl prostřední člen}]
\begin{enumerate}[label=\alph*)]
\item $2x^2+5x+2$;
\item $3x^2+10x+3$;
\item $2x^2-x-3$;
\item $4x^2-4x-3$.
\end{enumerate}
\tcblower
\begin{enumerate}[label=\alph*)]
\item $2x^2+4x+x+2=(2x+1)(x+2)$;
\item $3x^2+9x+x+3=(3x+1)(x+3)$;
\item $2x^2-3x+2x-3=(2x-3)(x+1)$;
\item $4x^2-6x+2x-3=(2x+1)(2x-3)$.
\end{enumerate}
\end{zadaniReseniBox}
\blok{Více metod v jediném rozkladu}
Některé výrazy nejprve upravíme známým vzorcem nebo vytýkáním a teprve potom seskupujeme. \textbf{Každý získaný činitel ještě zkontrolujeme}, zda jde rozložit dále.
\begin{prikladReseniBox}[breakable=false,title={Příklad 3: seskupování bez dalšího vzorce}]
Rozlož $x^3-4x^2+x-4$.
\tcblower
\begin{align*}
x^3-4x^2+x-4&=x^2(x-4)+(x-4)\\
&=(x^2+1)(x-4).
\end{align*}
Mnohočlen $x^2+1$ už nad reálnými čísly na lineární činitele nerozložíme.
\end{prikladReseniBox}
\begin{prikladReseniBox}[breakable=false,title={Příklad 4: poznej hotový vzorec v části výrazu}]
Rozlož $a^2+2ab+b^2-9$.
\tcblower
\begin{align*}
a^2+2ab+b^2-9&=(a+b)^2-3^2\\
&=(a+b-3)(a+b+3).
\end{align*}
Nejdříve rozpoznáme druhou mocninu, poté rozdíl čtverců.
\end{prikladReseniBox}
\begin{zadaniReseniBox}[breakable=false,title={Samostatná práce: vyber pořadí metod}]
\begin{enumerate}[label=\alph*)]
\item $x^3+2x^2-4x-8$;
\item $2a^2+4ab+2b^2-18$;
\item $x^4-5x^2+4$ (nápověda: $t=x^2$);
\item $x^5+x^3-x^2-1$ (rezerva).
\end{enumerate}
\tcblower
\begin{enumerate}[label=\alph*)]
\item $x^2(x+2)-4(x+2)=(x+2)(x-2)(x+2)=(x-2)(x+2)^2$;
\item $2[(a+b)^2-9]=2(a+b-3)(a+b+3)$;
\item $(x^2-1)(x^2-4)=(x-1)(x+1)(x-2)(x+2)$;
\item $x^3(x^2+1)-(x^2+1)=(x^2+1)(x^3-1)=(x^2+1)(x-1)(x^2+x+1)$.
\end{enumerate}
\end{zadaniReseniBox}
\blok{Volitelná rezerva: zkoušej i jiná seskupení}
\begin{zadaniReseniBox}[breakable=false,title={Rozšiřující příklad: složitější seskupování}]
Rozlož $x^3-3x+2$. Najdi rozdělení prostředního členu, které umožní seskupit.
\tcblower
Rozdělíme $-3x=-x-2x$:
\begin{align*}
x^3-3x+2&=x^3-x-2x+2\\
&=x(x^2-1)-2(x-1)\\
&=(x-1)[x(x+1)-2]\\
&=(x-1)(x^2+x-2)=(x-1)^2(x+2).
\end{align*}
\end{zadaniReseniBox}
\blok{Domácí procvičení}
\begin{zadaniReseniBox}[breakable=false,title={Domácí úkol: od rozdělení členu ke kombinaci metod}]
\begin{enumerate}[label=\alph*)]
\item $2x^2+9x+4$;
\item $3x^2+7x+2$;
\item $2x^2-7x+3$;
\item $6x^2+x-2$;
\item $x^3+5x^2-4x-20$;
\item $a^2-2ab+b^2-16$;
\item $x^4-10x^2+9$;
\item $x^3-2x^2-9x+18$.
\end{enumerate}
\tcblower
\begin{enumerate}[label=\alph*)]
\item $(2x+1)(x+4)$;
\item $(3x+1)(x+2)$;
\item $(2x-1)(x-3)$;
\item $(3x+2)(2x-1)$;
\item $x^2(x+5)-4(x+5)=(x+5)(x-2)(x+2)$;
\item $(a-b)^2-4^2=(a-b-4)(a-b+4)$;
\item $(x^2-1)(x^2-9)=(x-1)(x+1)(x-3)(x+3)$;
\item $x^2(x-2)-9(x-2)=(x-2)(x-3)(x+3)$.
\end{enumerate}
\end{zadaniReseniBox}
\end{document}
